\chapter{Background}\label{r03:ch:background}

\emph{The terms the rest of the report relies on: the lattice Boltzmann method and its two
ways of moving data, how waLBerla stores and exchanges the data, the software that generates
the kernel, the RISC-V vector extension, the memory system, and the measurement terms. Readers
who know a section's subject can skip it; the glossary (Appendix~\ref{r03:app:glossary}) repeats
every definition in one line.}

\section{The lattice Boltzmann method}\label{r03:sec:bg-lbm}

The lattice Boltzmann method (LBM) simulates a fluid on a regular grid of \term{cells}. Each
cell does not store the fluid's velocity directly. It stores $q$ numbers $f_i$, one per
\term{lattice direction} $c_i$. Each $f_i$ says how much fluid in the cell moves along
$c_i$. They are called \term{particle distribution functions}, abbreviated
\term{PDFs}. (In this report ``PDF'' always means such a value, never the file format.)

\paragraph{D3Q19.} The model used throughout is \term{D3Q19}: three dimensions, 19
directions per cell. The directions are the rest direction $(0,0,0)$, the 6 face
neighbours such as $(1,0,0)$, and the 12 edge neighbours such as $(1,1,0)$. Every direction
$i$ has an \term{inverse direction} $\bar\imath$ with $c_{\bar\imath} = -c_i$. Each
direction also has a \term{weight} $w_i$ ($1/3$ for the rest direction, $1/18$ for a face,
$1/36$ for an edge).

\paragraph{Macroscopic values.} The fluid's \term{density} $\rho$ and \term{velocity} $u$
in a cell follow from its PDFs:
\[
  \rho = \sum_i f_i, \qquad \rho\,u = \sum_i c_i f_i .
\]

\paragraph{A time step} applies two operations to every cell:
\begin{description}
  \item[Collision] moves the PDFs of a cell toward an \term{equilibrium} $f^{\mathrm{eq}}_i(\rho,
        u)$, a fixed polynomial in $u$. The textbook form is $f^{\mathrm{eq}}_i = w_i \rho\,(1 +
        3\,c_i\cdot u + \tfrac92 (c_i\cdot u)^2 - \tfrac32 u^2)$. With the
        \term{single-relaxation-time} method (\term{SRT}), all PDFs relax at the same rate
        $\omega$, which is set by the fluid's viscosity:
        \[
          f_i \leftarrow f_i + \omega\,\bigl(f^{\mathrm{eq}}_i - f_i\bigr).
        \]
        The \term{two-relaxation-time} method (\term{TRT}) uses one rate for the symmetric and
        another for the antisymmetric part of each pair $i$, $\bar\imath$. It appears only in
        the tests (\cref{r03:sec:wl-tests}).
  \item[Streaming] moves each post-collision value $f_i$ to the neighbouring cell
        $x + c_i$.
\end{description}
One update of all cells is one time step. Collision and streaming of all cells together are
done by one generated function, the \term{stream/collide kernel}. One pass of that kernel
over all cells is called a \term{sweep}.

\paragraph{Work per cell.} A D3Q19 cell update loads 19 PDFs, does about 200 floating-point
operations, and stores 19 PDFs. The benchmark also stores the density and the three velocity
components into separate output fields, so it makes 23 stores per cell. These density and
velocity fields are called the \term{output} below; the benchmark only reads them after the
last step, to compute its checksum (\cref{r03:sec:app}).

\paragraph{Zero-centred storage.} lbmpy stores $f_i - w_i$ instead of $f_i$, and
$\delta\rho = \rho - 1$ instead of $\rho$. This keeps the stored numbers small and improves
the rounding behaviour, especially in single precision. It does not change the method.

\section{Two ways to stream: pull and AA}\label{r03:sec:bg-stream}

Streaming moves data between cells. A kernel can do it in two ways, shown in
\cref{r03:fig:timestep}:
\begin{description}
  \item[Pull] uses two PDF arrays. For each cell $x$ and direction $i$, the kernel reads the
        value that streams into $x$, which is $f_i$ of the neighbour $x - c_i$. It collides
        the 19 values and writes the results into cell $x$ of the second array. Then the two
        arrays swap roles. Every value is read from one array and written to another.
  \item[AA pattern] \autocite{bailey-aa} uses one array, updated in place. It alternates two
        kinds of step, whose read and write positions are chosen so that each cell writes
        exactly the storage places it has read (\cref{r03:fig:aa}):
        \begin{itemize}
          \item \term{even step}: read $f_i$ from cell $x$, and write the post-collision value
                of direction $i$ into the place of $\bar\imath$, also in cell $x$;
          \item \term{odd step}: read $f_{\bar\imath}$ from cell $x - c_i$, and write the
                post-collision value of direction $i$ into the place of $i$ in cell
                $x + c_i$.
        \end{itemize}
        After an even and an odd step every value is back in its usual place and has moved
        by two cells.
\end{description}
The storage place of one direction's value in one cell is called a \term{slot}: ``slot $i$
of cell $x$''. AA streaming needs half the memory and, as \cref{r03:sec:lever-aa} shows,
fewer concurrent memory accesses.

\begin{figure}[htbp]
  \centering
  \resizebox{\linewidth}{!}{% The timed loop of SweepgenLbmCPU.cpp: pull (two PDF fields) and AA (one field, two kinds of step).
\begin{tikzpicture}[node distance=4mm and 5mm, every node/.style={font=\scriptsize}]
  % pull
  \node[font=\footnotesize\bfseries,anchor=east] (lp) at (-0.3,0) {pull};
  \node[r03box,text width=34mm] (p1) at (2.2,0)
    {\texttt{StreamCollide}\\reads \texttt{f}, writes \texttt{f\_tmp}, $\rho$, $u$};
  \node[r03box,text width=18mm,right=of p1] (p2) {swap \texttt{f}, \texttt{f\_tmp}};
  \node[r03box,text width=30mm,right=of p2] (p3) {\texttt{streamPullSync}\\halo exchange};
  \draw[r03arrow] (p1) -- (p2); \draw[r03arrow] (p2) -- (p3);
  \draw[r03arrow] (p3.east) -- ++(0.35,0) |- ++(0,0.75) -| (p1.north);
  \node[r03label,anchor=south] at ($(p1.north)!0.5!(p3.north)+(0,0.75)$) {every step};

  % AA
  \node[font=\footnotesize\bfseries,anchor=east] (la) at (-0.3,-2.0) {AA};
  \node[r03box,text width=34mm] (a1) at (2.2,-2.0)
    {\texttt{StreamCollideEven}\\slots of cell $x$, in place};
  \node[r03box,text width=24mm,right=of a1] (a2) {\texttt{aaPullSync}\\fills ghost layers};
  \node[r03box,text width=34mm,right=of a2] (a3)
    {\texttt{StreamCollideOdd}\\reads at $x - c_i$, writes at $x + c_i$};
  \node[r03box,text width=24mm,right=of a3] (a4) {\texttt{aaPushSync}\\returns ghost writes};
  \draw[r03arrow] (a1) -- (a2); \draw[r03arrow] (a2) -- (a3); \draw[r03arrow] (a3) -- (a4);
  \draw[r03arrow] (a4.south) -- ++(0,-0.4) -| (a1.south);
  \node[r03label,anchor=north] at ($(a1.south)!0.5!(a4.south)+(0,-0.4)$) {two steps};
\end{tikzpicture}
}
  \caption[One time step with pull and with AA streaming]{One time step of the benchmark with
    pull streaming (two PDF arrays \texttt{f} and \texttt{f\_tmp}, swapped after every step)
    and with AA streaming (one array, even and odd steps alternating). After every step, a
    halo exchange (\cref{r03:sec:bg-blocks}) brings the ghost layers up to date. The names
    in typewriter font are the generated kernels and the waLBerla functions involved.}\label{r03:fig:timestep}
\end{figure}

\begin{figure}[htbp]
  \centering
  \resizebox{\linewidth}{!}{% The AA pattern for the pair E (c = +x) / W (c = -x) around cell x, as in lbmpy's AA accessors and
% hand/hand-aa.hpp: even step reads slot i at x and writes slot inv(i) at x; odd step reads slot
% inv(i) at x - c_i and writes slot i at x + c_i.
\begin{tikzpicture}[every node/.style={font=\scriptsize},
    slot/.style={draw=psgray!80,fill=codebg,minimum width=8mm,minimum height=6mm,inner sep=0pt},
    coll/.style={circle,draw=psgray!90,fill=white,inner sep=1.5pt,align=center},
    rd/.style={-{Stealth[length=1.8mm]},thick,draw=psblue!85!black,dashed},
    wr/.style={-{Stealth[length=1.8mm]},thick,draw=psorange!90!black}]
  % three cells per panel, two slots per cell; prefixes e (even) and o (odd)
  \foreach \p/\px in {e/0, o/7.4} {
    \foreach \c/\lab in {0/$x-1$, 1/$x$, 2/$x+1$} {
      \node[slot] (\p W\c) at (\px+\c*2.1+0.4,0) {W};
      \node[slot] (\p E\c) at (\px+\c*2.1+1.2,0) {E};
      \node[r03label] at (\px+\c*2.1+0.8,-0.55) {cell \lab};
    }
    \node[coll] (\p col) at (\px+2.9,1.6) {collide\\at $x$};
  }
  \node[font=\footnotesize\bfseries] at (2.9,2.45) {even step: in place};
  \node[font=\footnotesize\bfseries] at (10.3,2.45) {odd step: from and to the neighbours};

  % even: read E and W at x, write post-collision E into slot W and W into slot E, both at x
  \draw[rd] (eE1.north) to[bend right=12] (ecol);
  \draw[rd] (eW1.north) to[bend left=12] (ecol);
  \draw[wr] (ecol.west) to[bend right=40] node[pos=0.35,left,text=psorange!90!black] {E} (eW1.north west);
  \draw[wr] (ecol.east) to[bend left=40] node[pos=0.35,right,text=psorange!90!black] {W} (eE1.north east);

  % odd: direction E is read from slot W at x-1 and written to slot E at x+1; W the other way
  \draw[rd] (oW0.north) to[bend left=8] (ocol);
  \draw[rd] (oE2.north) to[bend right=8] (ocol);
  \draw[wr] (ocol.east) to[bend left=25] node[pos=0.5,above,text=psorange!90!black] {E} (oE2.north west);
  \draw[wr] (ocol.west) to[bend right=25] node[pos=0.5,above,text=psorange!90!black] {W} (oW0.north east);

  % legend
  \draw[rd] (0.6,-1.2) -- ++(0.7,0) node[right,text=black] {read (pre-collision value)};
  \draw[wr] (5.4,-1.2) -- ++(0.7,0) node[right,text=black]
    {write (post-collision value of the direction named on the arrow)};
\end{tikzpicture}
}
  \caption[The AA pattern for one pair of directions]{The AA pattern for one pair of
    directions, east (E, $c = (1,0,0)$) and west (W, $c = (-1,0,0)$), seen from cell $x$.
    Each cell has one slot per direction; only the E and W slots are drawn. Even step: both
    values are read from cell $x$ and written back to cell $x$, each into the other's slot.
    Odd step: values are read from and written to the neighbours.}\label{r03:fig:aa}
\end{figure}

\section{How waLBerla stores and exchanges the data}\label{r03:sec:bg-blocks}

\paragraph{Fields.} waLBerla stores the PDFs of all cells in a \term{field}: a
four-dimensional array with the indices $x$, $y$, $z$ and the direction $f$. The benchmark
uses the \term{fzyx layout}: the direction index varies slowest and $x$ fastest. So the
values of one direction form one contiguous three-dimensional array, called a
\term{slice} below; a D3Q19 field has 19 slices. Inside a slice, cells that are neighbours
in $x$ are neighbours in memory. A \term{row} is the run of cells of one slice with fixed
$y$ and $z$. Rows are padded to a multiple of 16~bytes.

\paragraph{Blocks and ghost layers.} The simulation domain is cut into \term{blocks}, and
each process works on its own block. Each block's field has one extra layer of cells
around it, the \term{ghost layer}. Ghost cells hold copies of the values from the
neighbouring block, so that the kernel can read across the block border without special
cases.

\paragraph{Halo exchange.} After every time step the ghost layers must be refreshed: values
near the border of each block are copied into the ghost layer of its neighbour. This is the
\term{halo exchange}. Between processes it uses MPI (Message Passing Interface), the standard
library for sending messages between processes. In waLBerla, a \term{pack info} class
decides which values cross a block face, in three operations:
\begin{itemize}
  \item \code{pack}: copy the values to send into a message buffer;
  \item \code{unpack}: copy a received buffer into the ghost layer;
  \item \code{localCopy}: copy directly when sender and receiver are the same process.
\end{itemize}
The benchmark's domain is \term{periodic}: what leaves at one side comes back in at the
opposite side. With a single block, every face of the block is its own neighbour, so the
whole halo exchange consists of \code{localCopy} calls.

\section{The software that generates the kernel}\label{r03:sec:bg-stack}

\begin{description}
  \item[lbmpy] (Python) derives the LBM update of one cell symbolically, as a list of
        assignments, from a description of the method: lattice (D3Q19), collision (SRT or
        TRT), streaming (pull or AA), output fields \autocite{lbmpy}. Its option
        \term{global common-subexpression elimination} (\term{CSE}) computes every repeated
        sub-expression of the whole update once.
  \item[\ps{}] (Python) turns such assignments into a C++ function that loops over all cells
        of a field. It supports several \term{targets}: plain C++ for any
        CPU, x86 vector instructions, RVV (\cref{r03:sec:bg-rvv}) and GPUs. For vector targets
        it vectorizes the innermost loop itself, by writing vector intrinsics
        (\cref{r03:sec:bg-rvv}) into the code. This is called \term{explicit vectorization},
        as opposed to \term{auto-vectorization}, where the C++ compiler tries to vectorize
        plain loops on its own.
  \item[sweepgen] (Python, part of waLBerla) wraps generated kernels into C++ \term{sweep}
        classes that waLBerla can apply to each block. Its building block
        (\term{prefab}) \code{LbmBulk} generates the standard LBM sweeps of a method:
        \code{StreamCollide} (one time step) and \code{InitPdfs} (set the PDFs to the
        equilibrium of given density and velocity fields).
  \item[waLBerla] (C++) provides blocks, fields, the halo exchange and the time loop
        \autocite{walberla}. The benchmark uses waLBerla's \term{v8 core}, an experimental
        new core library in \file{include/walberla/v8/} that is enabled with the CMake option
        \code{WALBERLA_ENABLE_V8CORE}.
\end{description}
The path from physics to machine code is therefore: lbmpy (formulas) $\rightarrow$ \ps{}
(loop kernel in C++) $\rightarrow$ sweepgen (sweep class) $\rightarrow$ the C++ compiler
$\rightarrow$ the waLBerla application.

\section{RISC-V and its vector extension}\label{r03:sec:bg-rvv}

\paragraph{RISC-V} is an open instruction set. \code{rv64gc} names its general-purpose 64-bit
form with integer multiplication, atomics, single- and double-precision floating point, and
compressed instructions. A program compiled with \code{-march=rv64gc} uses no vector
instructions; this is how the plain scalar build is compiled. Further \term{extensions} add
instructions; a core lists the extensions it supports in its \term{ISA string}, shown in
\file{/proc/cpuinfo} on Linux.

\paragraph{The vector extension (RVV)} adds 32 vector registers. Their width in bits, the
\term{VLEN}, depends on the core; on the X60 it is 256. A register is divided into
\term{lanes} of the \term{element width}: 8 lanes of 32 bits (written \term{e32}, used for f32)
or 4 lanes of 64 bits (\term{e64}, for f64). One vector instruction operates on all lanes.
The \term{register grouping} \term{LMUL} lets one instruction operate on 2, 4 or 8 registers at
once, as one long register; LMUL~1 means one register.

\paragraph{Instructions used in this report.}
\begin{description}
  \item[\code{vle32}, \code{vse32}] load and store a register's worth of consecutive 32-bit
        values (\term{unit-stride} access); \code{vle64}/\code{vse64} for 64-bit values.
  \item[\code{vfadd}, \code{vfmul}, \code{vfdiv}] floating-point add, multiply, divide. The
        \term{.vv form} takes two vector operands; the \term{.vf form} takes a vector and a
        scalar floating-point register, whose value is used in every lane.
        \code{vfrsub}/\code{vfrdiv} are the .vf forms with the operands reversed
        ($s - v$, $s / v$).
  \item[\code{vfmv.v.f}] copies a scalar into every lane of a vector register (a
        \term{splat} or broadcast).
  \item[\code{vfmacc}, \code{vfmsac}, \code{vfnmsac}] \term{fused multiply-add} (\term{FMA}):
        $d \leftarrow d + a b$, $d \leftarrow a b - d$, $d \leftarrow d - a b$ in one
        instruction, rounded once instead of twice.
\end{description}
C++ code uses these instructions through \term{intrinsics}: functions that the compiler
provides (header \file{riscv_vector.h}) and translates to one instruction each
\autocite{rvv-intrinsics}. For example,
\code{__riscv_vle32_v_f32m1} loads 8 floats into a value of type \code{vfloat32m1_t}
(f32, LMUL~1). The compiler allocates registers and orders the instructions, but it treats
each intrinsic as a fixed operation: it does not merge a multiply and an add into an FMA by
itself.

\paragraph{Other extensions named in this report.}
\begin{description}
  \item[Zicbop] cache-block prefetch: \code{prefetch.r} and \code{prefetch.w} ask the cache to
        load a line for reading or writing. They are encoded as \term{HINTs}: instruction
        encodings that do nothing architecturally (here \code{ori x0, rs, imm}, which writes to
        the always-zero register). A core without Zicbop therefore executes them as no-ops.
  \item[Zicbom, Zicboz] cache-block management and zeroing; \code{cbo.zero} sets a whole
        cache line to zero without reading it from memory.
  \item[\code{zvl256b}] in a \code{-march} string: VLEN is at least 256 bits, so the compiler
        may assume 8 f32 lanes.
\end{description}

\paragraph{Scheduling model.} A compiler orders instructions using a description of the target
core's latencies and execution units, its \term{scheduling model}. The option
\code{-mtune=<core>} selects it. Clang~18 has a model for the X60
(\code{-mtune=spacemit-x60}); GCC~14 does not.

\section{The memory system}\label{r03:sec:bg-mem}

\begin{description}
  \item[Cache line] the unit in which data moves between memory and the caches: 64~bytes on
        the X60.
  \item[Caches] each X60 core has a 32~KiB level-1 data cache (\term{L1D}). Four cores share
        a 512~KiB level-2 cache (\term{L2}). There is no third level. An access to a line that
        is not in a cache is a \term{miss}; the line then comes from the next level or from
        main memory (\term{DRAM}). A DRAM access takes several hundred cycles
        (\term{latency}).
  \item[In-order core] the X60 issues instructions in program order. When an instruction needs
        the result of a load that missed, it waits, and so does everything after it. The core
        can keep only a few misses outstanding at once (\term{misses in flight}).
  \item[Stream] a sequence of accesses that walks through memory with a constant step. Each
        row-by-row pass through one slice of the PDF field is one stream; a kernel that reads
        19 slices and writes 23 arrays has 42 concurrent streams.
  \item[Hardware prefetcher] a unit in the cache that detects streams and loads their next
        lines before the program asks for them, so that the misses overlap. It can track only
        a limited number of streams at once.
  \item[Software prefetch] an instruction placed in the program that asks the cache to load a
        line that will be needed soon, without waiting for it (\code{prefetch.r},
        \code{prefetch.w}).
  \item[Write-allocate] a store to a line that is not in the cache first loads that line from
        memory. Writing a second array therefore costs memory reads as well.
  \item[Bandwidth] bytes moved per second (GB/s).
  \item[Pages, TLB, huge pages] programs use virtual addresses, which the core translates in
        units of \term{pages} (4~KiB). Recent translations are kept in the \term{TLB}
        (translation lookaside buffer). \term{Huge pages} of 2~MiB need fewer translations.
        Linux can provide them automatically, as \term{transparent huge pages} (\term{THP});
        in its \code{madvise} mode only for memory that asks for them.
\end{description}

\section{Measurement terms}\label{r03:sec:bg-measure}

\begin{description}
  \item[MLUP/s] million lattice-cell updates per second: cells $\times$ time steps / seconds /
        $10^6$.
  \item[Cycles per cell, instructions per cell] clock cycles or executed instructions divided
        by cells $\times$ time steps. At 1.6~GHz, cycles per cell $= 1600 / \text{MLUP/s}$.
  \item[IPC] instructions per cycle.
  \item[PMU, perf] the \term{performance monitoring unit} (PMU) counts hardware events such as
        cycles, instructions and cache misses. Linux's \code{perf} tool reads it:
        \code{perf stat} counts events, \code{perf record} samples where the program spends
        its time, \code{perf annotate} shows the samples per instruction. The X60 has
        vendor-specific events in addition, such as \code{l2_miss}; \cref{r03:tab:pmu1}
        explains the ones used.
  \item[Backend stall] a cycle in which the core issues no instruction because an operand or an
        execution resource is not ready.
  \item[Best of 3] each program is run three times, and the highest MLUP/s is reported.
  \item[Session] one unattended job on the board that measures a group of variants back to back
        (\cref{r03:sec:hw}). The board's absolute speed varies by a few percent from session to
        session, so speedups are always computed between values of the same session.
  \item[QEMU] an emulator that runs RISC-V programs on an x86 computer, including the vector
        extension. It checks correctness, not speed.
\end{description}
